\section{Introduction}
\label{sec:introduction}

The GFM framework defines its actor as a system that maximizes
$\volL$, the poset measure over the capability
space~\citep[Definition~6]{lasser2026gfm}.  The axiomatization
of~\citep[Proposition~1]{lasser2026poset} establishes $\volL$ as
a poset measure satisfying axioms M1--M6 and proves it
self-balancing~\citep[Theorem~1]{lasser2026poset}.  The analysis
of~\citep{lasser2026horizon} extends to multi-substrate
collectives, proving that full domination is anti-maximizing
under discounting~\citep[Proposition~6]{lasser2026horizon}.
Throughout, $\volL$ is treated as a faithful proxy for the
experiential optionality the framework is meant to protect.

But $\volL$ measures what the collective \emph{can} do, not what
it \emph{does}.  The B-to-C gap, named
in~\citep[Section~7]{lasser2026gfm} and flagged as a primary
open problem in the sequence, is the possibility that these
diverge.  Two failure modes illustrate:

\begin{enumerate}
\item \textbf{The Doll Problem}
  \citep[Section~7.1]{lasser2026gfm}: a collective possesses
  rich capabilities but exercises few.  $\volL$ is high but
  experiential value is zero.

\item \textbf{Wireheading}
  \citep[Discussion]{lasser2026exo}: a collective exercises
  capabilities in narrow self-reinforcing loops that serve the
  measurement system rather than experiential value.  High
  $\volL$ and apparent exercise, but the exercise is degenerate.
\end{enumerate}

Both are instances of Goodhart's
Law~\citep{strathern1997improving,manheim2019categorizing}: when
the measure ($\volL$) diverges from the target (experiential
optionality), optimizing the measure no longer serves the
target.

\subsection{Why direct measurement fails}

A natural first approach is to define a \emph{realized capability
volume} $\volR$ by restricting $\volL$ to an exercised
sub-poset---capabilities counterfactually load-bearing for
realized outputs---and monitor the ratio $\beta = \volR/\volL$.
This fails on two grounds:

\paragraph{Tractability concern.}
The counterfactual query (``which capabilities are
load-bearing for which realized outputs?'') requires global
knowledge of every realized output the collective produces and
of every causal pathway from each capability to those outputs.
Establishing a complexity lower bound is outside this paper's
scope; the qualitative observation that direct computation
faces high practical cost at collective scale, alongside the
privacy concern below, motivates the indirect channel.

\paragraph{Privacy failure.}
Computing realized-output sets requires observational access to
what the collective actually does.  This is the
total-surveillance scenario.  The framework's operational
commitments~\citep{lasser2026exo,lasser2026phase,lasser2026wamura}
rule out panopticon-style observation.

\subsection{What this paper does}

We observe that agents \emph{voluntarily disclose}
$\volR$-content through an existing, privacy-respecting channel:
\textbf{sacrifice}.  When an agent surrenders a benchmarkable
capability (whose $\volL$ contribution is known) in exchange for
an unbenchmarked bundle (whose $\volR$ contribution is unknown),
S1 (free choice with non-degrading alternatives) makes the trade
voluntary, and revealed preference establishes a valuation
comparison ($\Ui(Y) \geq \Ui(X)$); composing this with S0
(valuation-bounded-by-exercise), S2 (benchmark grounding), and
S4(a) (valuation additivity) under the genuine-exchange
conditions of Definition~\ref{def:sacrifice_event} yields a
lower bound $\volR^{[W]}(Y) \geq \Delta\volL(X)$ on the
unbenchmarked bundle's exercise contribution.  Aggregating
across events with pairwise poset-independent bundles produces
a constructive lower bound on $\volR$ on the observed subset
$U$, equivalently an upper bound on the subset-local B-to-C
gap $1 - \betaL(U,T)$ (with $\volL(U)$ in the denominator),
or equivalently on the global conservative gap
$1 - \betaL(P,T)$ when the denominator is $\volL(P)$ and
$\volRlower$ is zero-extended outside $U$
(Definition~\ref{def:btoc_ratio}).

This paper formalizes the sacrifice event
(Section~\ref{sec:sacrifice_model}), states the six assumptions
(S0--S5) under which the aggregate bound holds, proves the
central theorem (Section~\ref{sec:lower_bound}), analyzes the
two observation modalities (money and time,
Section~\ref{sec:two_channels}), establishes the structural
non-self-balancing finding for $\volR$
(Section~\ref{sec:axiom_inheritance}), and composes the channel
with the commitment protocol
of~\citep{lasser2026exo}
(Section~\ref{sec:commitment}).

\subsection{Scope: what this paper does not do}
\label{sec:not_do}

\begin{itemize}
\item
\textbf{It does not filter S1-failing below-sufficiency
sacrifices.}
Not every observed sacrifice event is voluntary---agents pay
below-sufficiency costs (dehydration, debt-servicing) when no
non-degrading alternative is available, in which case S1 fails
(retention is degrading and no third non-degrading option
exists).  Such events have reversed polarity relative to the
theorem's assumptions.  The companion
paper~\citep[Need~Sufficiency]{lasser2026paper8b} introduces the
\emph{need-sufficiency architecture} that supplies the
structural condition for diagnosing this S1 failure mode; the
filtering itself remains an institutional-attestation step.
Below-sufficiency targeting where a third non-degrading
alternative is available may still satisfy S1; classification
of such events likewise routes through institutional
attestation.  The present paper takes the filter as an
external precondition: events reaching the aggregate are
assumed S1-admissible.

\item
\textbf{It does not decompose the B-to-C gap into diagnostic
categories.}  Interpreting why $\betaL = \volRlower/\volL$ is
low requires partitioning capabilities into restricted,
covered, dormant, and residual cells and tracking which cell
dominates the gap.  The gap decomposition machinery and its
wireheading-consistent concentration signal are developed
in~\citep[Need~Sufficiency]{lasser2026paper8b}.

\item
\textbf{It does not provide a complete theory of experiential
value.}  The residual class (capabilities structurally outside
the sacrifice channel's reach) is characterized
in~\citep[Need~Sufficiency]{lasser2026paper8b}.

\item
\textbf{It does not replace $\volL$ with $\volR$.}  The framework
continues to optimize $\volL$; the revealed-sacrifice lower
bound is a diagnostic that detects when $\volL$ optimization has
gone wrong.  $\volR$ is the audit, not the objective.

\item
\textbf{It does not measure $\volR$ directly.}  The paper
produces a \emph{lower bound} on $\volR$ from trade events, not
a precise measurement.  An \emph{exact} per-capability exercise
witness $\volRexact$ within consented perimeters, together
with the boundary-closure conditions and OI--exercise bridge
that license combining $\volRexact$ with the bound's
$\volRlower$ outside the perimeter into a unified diagnostic
$\volR^{\mathrm{combined}}$, is developed in
Appendix~\ref{app:exercise_indicator}.  The combined diagnostic
is the natural deployment object for governance frameworks that
operate both inside and outside consented perimeters; main-text
sections~\ref{sec:lower_bound}--\ref{sec:axiom_inheritance}
state results for the privacy-respecting lower-bound channel
alone, and forward-point to the appendix where the
combined-diagnostic infrastructure is needed.
\end{itemize}

\subsection{Summary of results}

\begin{enumerate}
\item \textbf{Privacy-respecting observation channel with
  formal-plus-institutional privacy split}
  (Section~\ref{sec:sacrifice_model}).
  Revealed-sacrifice observation defines the sacrifice event and
  states six assumptions (S0--S5).  Three formal privacy
  properties are established (Theorem~\ref{thm:privacy}: no
  interior access, emission-rate discretization---bounded by
  trade frequency rather than observation rate, not a
  per-event information-leakage bound, commitment-layer
  composability); two further properties hold operationally
  only
  under named institutional preconditions
  (Proposition~\ref{prop:privacy_institutional}: consent via
  participation, third-party observability).  The
  ``privacy-minimal'' interpretation is conditional on the
  institutional preconditions; the formal layer alone establishes
  P1, P3, P4.

\item \textbf{Aggregate B-to-C Lower Bound}
  (Section~\ref{sec:lower_bound}).
  Theorem~\ref{thm:aggregate_bound} lower-bounds
  $\volR^{U,[W]}$---the window-active realized capability
  volume on the subset $U$ covered by S1-admissible trade
  events (free-choice voluntary trades with at least one
  non-degrading alternative in the choice set) with
  \emph{pairwise poset-independent bundles}---by aggregated
  sacrifice magnitudes.  Via
  Definition~\ref{def:btoc_ratio}, the bound induces two
  associated ratios: the subset-local exercised fraction
  $\betaL(U, T) = \volRlower(U, T) / \volL(U)$ on the observed
  $U$, and the global conservative fraction $\betaL(P, T) =
  \volRlower(U, T) / \volL(P)$ on the full capability-space
  denominator.  The subset-local ratio is the lower bound
  restricted to $U$; the global ratio is a lower bound on the
  true whole-system exercised fraction.  Both are lower bounds,
  not tight estimates: no sharpness/optimality theorem is
  claimed.  We write simply $\betaL$ when the choice of
  denominator is fixed by context.  Part~A
  (bundle-level, requires S0--S2 and S4(a)) is the main result;
  Part~B (per-capability, requires S0--S5 plus within-bundle
  poset-independence) is a conditional refinement.  The
  constructed lower bound is non-decreasing in observed
  admissible events: Part~A via sub-batch supremum
  (Proposition~\ref{prop:monotone_part_a}) and Part~B via
  per-capability max-attribution
  (Proposition~\ref{prop:monotone}).

\item \textbf{Two sacrifice channels, substitutable with
  residuals}
  (Section~\ref{sec:two_channels}).
  Money-sacrifice and time-sacrifice are substitutable via
  cooperative markets (employment converts time to money;
  hiring converts money to time).  The channels triangulate the
  same underlying substrate rather than covering disjoint
  capability spaces.  Non-fungibility residuals
  (agent-specific time, capital-concentrated access) are
  existential witnesses enabling strict improvement of the
  joint aggregate conditional on observed-coverage content
  (Corollary~\ref{cor:joint_coverage}).

\item \textbf{Non-self-balancing finding}
  (Section~\ref{sec:axiom_inheritance}).
  $\volR$ is non-invariant under capability removal:
  unexercised capabilities can be eliminated without reducing
  $\volR$, while $\volL$ contracts by their weight.  The
  self-balancing property therefore does not transfer from
  $\volL$ to $\volR$, so $\volR$ is a diagnostic, not an
  objective.  This structural finding is independent of the
  observation mechanism.

\item \textbf{Commitment-layer composition}
  (Section~\ref{sec:commitment}).
  The sacrifice channel composes with the commitment protocol
  of~\citep{lasser2026exo} to produce a privacy-minimal
  aggregate on distinct-category batches
  (Propositions~\ref{prop:commit_preserves}
  and~\ref{prop:zk_aggregation}).  The rollup ZK-certifies the
  Part~A formula inputs; the theorem's non-ZK hypotheses
  require external attestation
  (Remark~\ref{rem:external_attestation}).
\end{enumerate}

\subsection{Design principle}
\label{sec:design_principle}

This paper shares a structural discipline
with~\citep{lasser2026phase} and~\citep{lasser2026wamura}:
\textbf{resolution through frequency, not through depth.}
Controlled relaxation~\citep{lasser2026wamura} gains precision
by running more tests, not by surveilling each test more
intrusively.  Revealed-sacrifice observation gains precision by
observing more trades, not by observing each trader more
intrusively.

\subsection{Dependencies on prior papers}
\label{sec:dependencies}

A table summarising the prior-paper results this paper relies on
is collected in Appendix~\ref{app:notation}.  The companion
paper~\citep[Need~Sufficiency]{lasser2026paper8b} and the earlier sequence
papers~\citep{lasser2026gfm,lasser2026poset,lasser2026horizon,lasser2026scm,lasser2026exo,lasser2026phase,lasser2026wamura}
are presently distributed as preprints.  The propositions and
definitions imported here are pinned by their current
formulations; if the upstream preprints revise their numbering
or content, downstream references will need refreshing
correspondingly.  Stable archival identifiers will be added
as the sequence reaches archival status.
